% ===================================================================
%  TABELUL Nr. 11 — Orașul și clădirile lui. — Incendiul.
%  Traduction 2026 d'apres texto/io, controlee sur texto/fr, texto/en
%  et texto/es. Consigne : voir texto/ro/10-tabelul-01.tex.
%
%  OUVERTURE DE LA TROISIEME SERIE : « A TREIA SERIE ».
%
%  LE BLOC c1-03-1 EST CELUI QUI A COUTE UNE INVERSION AU FINNOIS ET
%  UNE AU BASQUE. Le roumain postpose son genitif comme l'ido —
%  « poarta (8) vechiului castel intarit (7) » — et ne paie rien.
% ===================================================================

%%K t11-apar-1 apar
\VUsaut{2.0mm}
\VUcentre{13.2pt}{90}{A TREIA SERIE}
\VUsaut{3.0mm}
\VUfilet{16mm}
\VUsaut{5.6mm}
\VUtitre{15.0pt}{60}{TABELUL Nr. 11}
\VUsaut{1.4mm}
\VUpk{11.4pt}{40}{Orașul și clădirile lui. --- Incendiul.}
\VUsaut{2.2mm}

%%K t11-tit-1 sub
\VUpk{10.2pt}{40}{Marginile.}
\VUsaut{3.0mm}

%%K t11-c1-01-1 p
1. --- Locuiesc într-un oraș mare la o sută de kilometri de ocean, care totuși este
un \VUgras{port} de mâna întâi \textsuperscript{(1)}. Râul care curge de-a lungul
lui are patru sute de metri lățime și face o radă frumoasă.

%%K t11-c1-02-1 p
2. --- Toată partea orașului de-a lungul \VUgras{cheiurilor}
\textsuperscript{(2)} are o înfățișare cu totul marinărească și industrială și
alcătuiește o \VUgras{mahala} \textsuperscript{(3)}, unde locuiesc numai marinari
și lucrători. Aceștia din urmă lucrează la \VUgras{uzine} \textsuperscript{(4)} și
la \VUgras{fabrica de gaz} \textsuperscript{(5)}, al cărei \VUgras{coș} înalt
\textsuperscript{(6)} și al cărei \VUgras{gazometru} se văd de departe.

%%K t11-c1-03-1 p
3. --- În Evul Mediu orașul era întărit, și se mai găsesc câteva rămășițe ale
metereselor sale, mai ales \VUgras{poarta} \textsuperscript{(8)} vechiului
\VUgras{castel întărit} \textsuperscript{(7)} cu două \VUgras{turnuri}
\textsuperscript{(9)} în laturi, care acum slujesc de \VUgras{închisoare}.

%%K t11-c1-04-1 p
4. --- În tot acest \VUgras{cartier}, la vedere foarte vechi, o singură clădire este
mai nouă: \VUgras{vama} \textsuperscript{(10)}. Stă între chei și o \VUgras{stradă}
mică \textsuperscript{(11)}, care duce la vechea \VUgras{primărie}
\textsuperscript{(12)}. Această din urmă clădire se recunoaște lesne după
\VUgras{turnul de veghe} uriaș \textsuperscript{(13)}, care se înalță deasupra
orașului vechi.

%%K t11-c1-05-1 p
5. --- De cealaltă parte a străzii \VUgras{stă} o clădire în același stil ca vama:
este \VUgras{bursa} \textsuperscript{(14)}. Una dintre fațadele ei dă spre o piață
mică, unde se află \VUgras{prefectura} \textsuperscript{(15)}. Orașul nostru, ce e
drept, nu este capitală, dar este totuși un însemnat centru de județ.

%%K t11-tit-2 sub
\VUsaut{5.0mm}
\VUpk{10.2pt}{40}{Partea de sus a orașului.}
\VUsaut{3.4mm}

%%K t11-c1-06-1 p
6. --- Garnizoana, destul de mare, este așezată în \VUgras{cazărmi} întinse și
foarte sănătoase \textsuperscript{(16)}; se întind până la gratiile
\VUgras{parcului orășenesc} \textsuperscript{(17)}. \VUgras{Bulevardul}
\textsuperscript{(19)}, care duce spre acest parc, trece pe în dosul \VUgras{casei
de economii} \textsuperscript{(18)} și iese în piața \VUgras{catedralei}
\textsuperscript{(20)}.

%%K t11-c1-07-1 p
7. --- Toate aceste clădiri se ridică în trepte pe un deal mic, unde stă și marele
nostru \VUgras{liceu} \textsuperscript{(21)}.

%%K t11-c1-07-2 p
Dacă cobori pe strada ușor pieptișă care iese în Strada Mare, ajungi la
\VUgras{tribunal} \textsuperscript{(22)}, înaintea căruia se întinde un
\VUgras{scuar} plăcut \textsuperscript{(26)}. Această \VUgras{piață} nu este mare,
dar în mijlocul orașului face o oază de verdeață și de răcoare. De aceea este
cercetată bucuros de orășeni, cărora le place să șadă sub șopronul
\VUgras{cafenelei} \textsuperscript{(25)}, ascultând fanfara militară care cântă
joia și duminica în \VUgras{chioșcul} \textsuperscript{(27)} dintre peluze și
straturi de flori.

%%K t11-c1-08-1 p
8. --- Pe una dintre aceste peluze stă o \VUgras{statuie} de bronz
\textsuperscript{(28)}, monument ridicat în amintirea unuia dintre orășenii noștri,
care a adus mari servicii orașului său de baștină.

%%K t11-tit-3 sub
\VUsaut{4.0mm}
\VUpk{10.2pt}{40}{Umblatul.}
\VUsaut{3.4mm}

%%K t11-c1-09-1 p
9. --- Orașul nostru nu este încă luminat cu electricitate; are numai
\VUgras{felinare cu gaz} \textsuperscript{(29)}. Dar \VUgras{presa locală} se
străduiește ca acest fel de luminat să fie îmbunătățit, \VUgras{așa cum} a fost
îmbunătățită \VUgras{tracțiunea} la tramvaie. \VUgras{Vechile} vagoane,
\VUgras{pe care le trăgeau} \VUgras{caii}, au fost înlocuite cu
\VUgras{tramvaie electrice} lesnicioase \textsuperscript{(31)}, care duc repede
călătorii dintr-un capăt \VUgras{al orașului în celălalt} pentru o plată foarte
\VUgras{mică}.

%%K t11-c1-10-1 p
10. --- Lângă tribunal este \VUgras{banca} \textsuperscript{(32)}, unde se
scontează polițele de negoț. \VUgras{Încasatorii băncii} \textsuperscript{(33)}
umblă din casă în casă ca să strângă ce se cuvine.

%%K t11-tit-4 sub
\VUsaut{4.0mm}
\VUpk{10.2pt}{40}{Arta și știința.}
\VUsaut{3.4mm}

%%K t11-c1-11-1 p
11. --- Clădirea frumoasă de alături este \VUgras{muzeul}. În ea se află deodată
\VUgras{biblioteca} orășenească \textsuperscript{(34)}, frumoasele
\VUgras{colecții de științe naturale} \textsuperscript{(35)} (muzeul de istorie
naturală) și eleganta \VUgras{pinacotecă} \textsuperscript{(36)}, care alcătuiește
o minunată \VUgras{expoziție} statornică.

%%K t11-c1-11-2 p
Este deschis publicului de două ori pe săptămână, dar vizitatorii îl pot cerceta
în orice zi pentru un mic bacșiș dat \VUgras{paznicului} \textsuperscript{(37)}.
\VUgras{Pictorii} \textsuperscript{(38)} vin acolo să lucreze. Îi vezi
\VUgras{înaintea} șevaletelor, cu \VUgras{paleta} în mână, copiind tablouri
vestite.

%%K t11-c1-12-1 p
12. --- Pe înălțimile din spatele muzeului se deosebesc abia \VUgras{cupola}
\textsuperscript{(40)} \VUgras{spitalului} \textsuperscript{(39)} și clădirile
\VUgras{școlii normale} \textsuperscript{(41)}, unde au învățat profesorii școlilor
noastre orășenești. În Piața Teatrului este \VUgras{poșta}
\textsuperscript{(43)}, așezată într-o \VUgras{casă particulară}
\textsuperscript{(42)}.

%%K t11-tit-5 sub
\VUsaut{5.0mm}
\VUpk{11.4pt}{40}{Incendiul.}
\VUsaut{3.0mm}

%%K t11-tit-6 sub
\VUpk{10.2pt}{40}{Strigătul «Foc!».}
\VUsaut{3.4mm}

%%K t11-c2-01-1 p
1. --- Prin oraș se răspândește vestea neliniștitoare: la teatru a izbucnit un
incendiu! Când ajung în \VUgras{piață} \textsuperscript{(96)}, mulțimea s-a strâns
deja pe \VUgras{trotuar} \textsuperscript{(46)} și pe \VUgras{caldarâm}
\textsuperscript{(47)}. \VUgras{Funcționarii poștali} \textsuperscript{(44)} și
\VUgras{poștașii} \textsuperscript{(45)} ies de la poștă. \VUgras{Vatmanul}
\textsuperscript{(48)} a fost silit să își oprească tramvaiul, ca să lase să treacă
\VUgras{salvarea} orășenească \textsuperscript{(50)}, venită cu
\VUgras{chirurgul} \textsuperscript{(52)} și cu \VUgras{infirmiera}
\textsuperscript{(51)}, ca să îl ajute pe \VUgras{rănit} \textsuperscript{(53)},
care este dus pe \VUgras{targă} \textsuperscript{(54)} spre un loc de lângă
\VUgras{chioșcul de ziare} \textsuperscript{(55)}.

%%K t11-c2-02-1 p
2. --- O companie de infanterie, care se întorcea de la instrucție, s-a oprit ca să
ajute la păstrarea rânduielii împreună cu \VUgras{gardiștii orășenești}
\textsuperscript{(61)} călări. \VUgras{Călăreții}, ai căror cai se ridică în două
picioare, se descurcă mai bine decât \VUgras{soldații} \textsuperscript{(57)} și
decât \VUgras{jandarmii} \textsuperscript{(63)}, împingând înapoi
\VUgras{gură-cască} \textsuperscript{(56)}. Jandarmii sunt pe deasupra foarte
ocupați. Unul dintre ei le arată \VUgras{polițiștilor} \textsuperscript{(64)} unde
să îl ducă pe celălalt \VUgras{păgubit} \textsuperscript{(65)}, un pompier viteaz
căzut de pe scară.

%%K t11-c2-03-1 p
3. --- \VUgras{Ofițerul} \textsuperscript{(66)} arată întocmai unde să fie așezată
\VUgras{pompa cu aburi} \textsuperscript{(67)}, care tocmai sosește. Așteptând
această mașină puternică, a poruncit să fie pusă \VUgras{pompa de mână}
\textsuperscript{(68)}, al cărei \VUgras{furtun} lung \textsuperscript{(69)} varsă
peste foc șuvoaie de apă. O pompează șase pompieri, iar tovarășul lor, ținând
\VUgras{țeava} \textsuperscript{(70)}, îndreaptă \VUgras{șuvoiul}
\textsuperscript{(71)} chiar în inima incendiului.

%%K t11-c2-04-1 p
4. --- De cealaltă parte a teatrului au ridicat o \VUgras{scară} mare
\textsuperscript{(72)}. Ea îngăduie pompierilor să se apropie de flacăra care se
înalță printre rotocoale de \VUgras{fum} negru \textsuperscript{(75)}.

%%K t11-tit-7 sub
\VUsaut{5.0mm}
\VUpk{10.2pt}{40}{Fapte îndrăznețe.}
\VUsaut{3.4mm}

%%K t11-c2-05-1 p
5. --- \VUgras{Focul} \textsuperscript{(74)} a început în foaierul
\VUgras{teatrului} \textsuperscript{(73)}, în vremea repetiției unei
\VUgras{opere}, a cărei primă \VUgras{reprezentație} trebuia să fie mâine. Din
fericire, în sală erau numai câțiva \VUgras{spectatori} \textsuperscript{(76)}.
Bieții oameni s-au adăpostit pe \VUgras{balcon} \textsuperscript{(77)}, unde vin
în ajutorul lor \VUgras{pompierii} viteji \textsuperscript{(86)} cu scara și cu
frânghii.

%%K t11-c2-06-1 p
6. --- Sub \VUgras{portic} \textsuperscript{(78)}, unde \VUgras{afișul}
\textsuperscript{(79)} vestește reprezentația, se vede cum fug \VUgras{muzicanții}
\textsuperscript{(81)} \VUgras{orchestrei} \textsuperscript{(80)}, ducându-și
instrumentele: \VUgras{vioara} \textsuperscript{(81)}, \VUgras{flautul}
\textsuperscript{(82)}, \VUgras{violoncelul} \textsuperscript{(83)},
\VUgras{trombonul} \textsuperscript{(84)}, \VUgras{toba} \textsuperscript{(85)} și
așa mai departe. Se încearcă să fie salvată o parte din \VUgras{mobilier}
\textsuperscript{(87)}.

%%K t11-c2-07-1 p
7. --- \VUgras{Actorii} \textsuperscript{(89)} și \VUgras{actrițele}
\textsuperscript{(88)}, \VUgras{cântăreții} și cântărețele au trebuit să iasă în
fugă în \VUgras{costumele} de scenă \textsuperscript{(90)}. Tenorul îi lămurește
\VUgras{ziaristului} \textsuperscript{(92)} cum s-a întâmplat.
\VUgras{Reporterul} a venit în fugă din cea dintâi clipă împreună cu autoritățile:
\VUgras{prefectul} \textsuperscript{(91)}, \VUgras{primarul}
\textsuperscript{(93)}, \VUgras{șeful poliției} \textsuperscript{(94)} și
\VUgras{judecătorul de instrucție} \textsuperscript{(95)}, care va face cercetarea
ca să afle cine poartă răspunderea.
\VUsaut{6.0mm}
\VUfilet{22mm}
